% The Baby Steps Camino organizers' manual.
% Build: xelatex manual.tex  (twice, for the contents page)
\documentclass[12pt]{article}

\usepackage{fontspec}
\setmainfont{DejaVu Sans}[Scale=1.02]
\newfontfamily\screenfont{DejaVu Sans}[Scale=0.98]
\usepackage[letterpaper,top=1in,bottom=1in,left=1.25in,right=1.25in]{geometry}
\usepackage{xcolor}
\usepackage{graphicx}
\usepackage{enumitem}
\usepackage{array}
\usepackage{tabularx}
\usepackage{booktabs}
% table columns that do not stretch their spaces
\newcolumntype{L}[1]{>{\raggedright\arraybackslash}p{#1}}
\newcolumntype{Y}{>{\raggedright\arraybackslash}X}
\usepackage{hyperref}

\definecolor{ink}{HTML}{1A1A1A}
\definecolor{accent}{HTML}{0B5F91}
\definecolor{softbg}{HTML}{EEF4F8}
\definecolor{warnbg}{HTML}{FDF1E3}
\definecolor{rule}{HTML}{C8D4DC}
\color{ink}
\hypersetup{colorlinks=true,linkcolor=accent,urlcolor=accent,pdftitle={Baby Steps Camino: the organizers' manual},pdfauthor={Baby Steps Camino}}

% generous leading, no paragraph indents, space between paragraphs instead
\linespread{1.22}\selectfont
\setlength{\parindent}{0pt}
\setlength{\parskip}{9pt}
\raggedright
\raggedbottom

% headings, without titlesec
\makeatletter
\renewcommand\section{\@startsection{section}{1}{0pt}%
  {28pt plus 4pt minus 2pt}{10pt}%
  {\normalfont\Large\bfseries\color{accent}}}
\renewcommand\subsection{\@startsection{subsection}{2}{0pt}%
  {20pt plus 3pt minus 2pt}{6pt}%
  {\normalfont\large\bfseries\color{ink}}}
\renewcommand\subsubsection{\@startsection{subsubsection}{3}{0pt}%
  {14pt plus 2pt minus 1pt}{4pt}%
  {\normalfont\normalsize\bfseries\color{ink}}}
\makeatother

% a boxed note. The optional argument is the background colour, the
% required one the heading word. Built with a save box, because a
% \colorbox argument cannot be split across \begin and \end.
\newsavebox{\notebox}
\newcommand{\notecolour}{softbg}
\newenvironment{note}[2][softbg]
  {\renewcommand{\notecolour}{#1}%
   \par\medskip
   \setlength{\fboxsep}{11pt}%
   \begin{lrbox}{\notebox}%
   \begin{minipage}{\dimexpr\linewidth-2\fboxsep\relax}%
   \setlength{\parskip}{7pt}%
   \textbf{#2}\par}
  {\end{minipage}%
   \end{lrbox}%
   \noindent\expandafter\colorbox\expandafter{\notecolour}{\usebox{\notebox}}%
   \par\medskip}

% a screenshot with a line of explanation under it
% Screenshots run wider than the text, into the margins, because the
% program's own type is small and a manual is no use if it cannot be read.
\newlength{\shotwidth}
\setlength{\shotwidth}{\dimexpr\textwidth+1.2in\relax}
\newcommand{\shot}[2]{%
  \par\medskip\noindent
  \makebox[\linewidth][c]{%
    \begingroup\setlength{\fboxsep}{2pt}%
    \fcolorbox{rule}{white}{\includegraphics[width=\dimexpr\shotwidth-2\fboxsep-2\fboxrule\relax]{#1}}%
    \endgroup}%
  \par\nopagebreak\smallskip\nopagebreak
  \makebox[\linewidth][c]{\begin{minipage}{\shotwidth}\small\itshape #2\end{minipage}}%
  \par\medskip}
% a close-up of one control, at its own natural size
\newcommand{\closeup}[3][\shotwidth]{%
  \par\medskip\noindent
  \makebox[\linewidth][c]{%
    \begingroup\setlength{\fboxsep}{2pt}%
    \fcolorbox{rule}{white}{\includegraphics[width=\dimexpr#1-2\fboxsep-2\fboxrule\relax]{#2}}%
    \endgroup}%
  \par\nopagebreak\smallskip\nopagebreak
  \makebox[\linewidth][c]{\begin{minipage}{\shotwidth}\small\itshape #3\end{minipage}}%
  \par\medskip}

% two phone screens side by side
\newcommand{\phonepair}[4]{%
  \par\medskip\noindent
  \makebox[\linewidth][c]{%
    \begingroup\setlength{\fboxsep}{2pt}%
    \fcolorbox{rule}{white}{\includegraphics[width=2.7in]{#1}}\hspace{0.45in}%
    \fcolorbox{rule}{white}{\includegraphics[width=2.7in]{#3}}%
    \endgroup}%
  \par\nopagebreak\smallskip\nopagebreak
  \makebox[\linewidth][c]{\begin{minipage}{2.7in}\small\itshape #2\end{minipage}%
    \hspace{0.45in}\begin{minipage}{2.7in}\small\itshape #4\end{minipage}}%
  \par\medskip}

% a word as it appears on screen
\newcommand{\btn}[1]{\textbf{#1}}
\newcommand{\siteroot}{https://forgave-infect.nisfeb.com}
\newcommand{\site}{\siteroot}

\begin{document}

% ---------------------------------------------------------------- cover
\thispagestyle{empty}
\vspace*{1.1in}
{\fontsize{30}{36}\selectfont\bfseries\color{accent} Baby Steps Camino\par}
\vspace{10pt}
{\fontsize{22}{28}\selectfont The organizers' manual\par}
\vspace{28pt}
{\large How to run registration, the wait list, the money and the three
days of the walk, using the Register program.\par}
\vfill
{\large
Everything in this manual is done from a web page. You need no special
software, only a computer or a phone and the address on the next page.\par}
\vspace{18pt}
{\small Written for the 2026 walk, December 4th, 5th and 6th.\par}
\newpage

% ---------------------------------------------------------------- contents
\pagestyle{empty}
{\Large\bfseries\color{accent} What is in this manual\par}
\vspace{6pt}
\begingroup
\setlength{\parskip}{0pt}
\renewcommand{\contentsname}{}
\tableofcontents
\endgroup
\newpage

\pagestyle{plain}
\setcounter{page}{1}

% ---------------------------------------------------------------- 1
\section{Start here}

Register is the program that takes sign-ups for the Baby Steps Camino
and helps you run the weekend. It does five jobs.

\begin{enumerate}[leftmargin=*,itemsep=6pt]
\item It takes sign-ups from pilgrims on a public web page.
\item It keeps the list of who is coming, and holds the wait list when a
  track fills up.
\item It keeps track of who has paid and who has signed the waiver.
\item It gives you counts, reports and spreadsheets.
\item On the three days of the walk it checks people in.
\end{enumerate}

\subsection{Who does what}

\begin{tabularx}{\linewidth}{@{}L{1.5in}Y@{}}
\toprule
\textbf{A pilgrim} & Signs up on the public page. Needs no account and
no password. Gets an email with a private link, and can change or
cancel their own registration from it. \\[6pt]
\textbf{You, an organizer} & Log in and use the backoffice: the roster,
the money, the reports, the settings. \\[6pt]
\textbf{A check-in volunteer} & On the three days, uses the check-in
app on a phone to tick people off and hand out wristbands. \\
\bottomrule
\end{tabularx}

\closeup[2.7in]{img/ov-pilgrim-landing.png}{What a pilgrim meets: the
public page, on a telephone.}

\subsection{The two addresses}

There are two web pages, and it matters which one you are on.

\begin{tabularx}{\linewidth}{@{}L{1.75in}Y@{}}
\toprule
\textbf{The public page} & \texttt{forgave-infect.nisfeb.com/apps/register/}
\newline This is what pilgrims see. Anybody may open it. \\[8pt]
\textbf{Your backoffice} & \texttt{forgave-infect.nisfeb.com/apps/register/admin}
\newline This is yours. It asks for a password. \\
\bottomrule
\end{tabularx}

Put the second address in your bookmarks. You will use it most days.

\begin{note}{Good to know}
If your addresses are different from the ones printed here, whoever set
the program up will have given you the right ones. Everything else in
this manual works the same way.
\end{note}

% ---------------------------------------------------------------- 2
\section{Logging in, and putting your name in}

\subsection{The password}

All the organizers share one password. It is a string of four short
words joined by dashes, like \texttt{lopdyt-monbyn-fabnyr-hasrev}. Ask
whoever set the program up for it, and keep it somewhere safe.

\begin{enumerate}[leftmargin=*,itemsep=6pt]
\item Open \texttt{forgave-infect.nisfeb.com/apps/register/admin}
\item If a login page appears, type the password and press the button.
\item The roster appears. You stay logged in on that computer for about
  a week, so you will not be asked again every day.
\end{enumerate}

\subsection{Your name}

Because you all share one password, the program cannot tell which of
you is which. So it asks you, once, on each computer or phone you use.

Look at the top right of the screen. If it says \btn{Set your name},
press it and type your first name.

\closeup{img/cu-name.png}{The name box. It appears by itself the first
time you try to change something.}

From then on the top right says \btn{Acting as Susan}, and every change
you make is signed with your name in that registration's history. If
someone later asks who cancelled the Brennans, the answer is on the
record.

\begin{note}{Careful}
The program remembers your name in that one web browser only. On a
different computer, or on your phone, it will ask again. That is
normal. Press the name at the top right any time to change it.
\end{note}

% ---------------------------------------------------------------- 3
\section{The Roster: your everyday screen}

The roster is the list of everyone who has signed up. It is the first
thing you see, and most days it is the only screen you need.

\shot{img/roster.png}{The whole roster screen. The names are links:
press one to open that registration. The parts of it are shown close up
on the next few pages.}

\subsection{The row of buttons along the top}

\closeup{img/cu-nav.png}{The row of buttons. It is the same on every
screen.}

Press any of them to change screens. \btn{Register} at the far left
always brings you back to the roster.

The small round dot beside your name is green when the page is
listening for changes. If another organizer changes something while you
are looking, your screen updates by itself within a few seconds. You
never need to reload the page.

\subsection{Segment: showing only part of the list}

\closeup{img/cu-filter.png}{The three boxes that narrow the list, and
the line of totals underneath.}

The box marked \btn{Segment} narrows the list. It starts on
\btn{registrations}, which means everyone who actually signed up: it
leaves out abandoned sign-ups and cancellations.

\begin{tabularx}{\linewidth}{@{}L{1.65in}Y@{}}
\toprule
\textbf{registrations} & Everyone real. Where you start. \\
\textbf{everything} & Adds abandoned sign-ups and cancellations. \\
\textbf{complete} & Paid and signed. Nothing left to do. \\
\textbf{pending} & Still owes a waiver, a payment, or an assistance
decision. \\
\textbf{wait list} & Waiting for a spot. \\
\textbf{financial assistance} & Asked for help with the fee. \\
\textbf{unpaid} & Has not paid. \\
\textbf{unsigned} & Has not signed the waiver. \\
\textbf{drafts} & Started the form and never finished. Worth an email. \\
\textbf{cancelled} & Cancelled. \\
\textbf{Bambino only} & The Sunday-only track. \\
\textbf{non-walkers} & Coming, but not walking. \\
\textbf{NE Florida Order of Malta and volunteers} & Local members and
volunteers. \\
\textbf{exempt} & Marked as not counting against the caps. \\
\textbf{admin adds} & Ones you added by hand. \\
\bottomrule
\end{tabularx}

\subsection{Search}

Type a name, an email address, a phone number or a parish into the
\btn{Search} box. It looks for every word you type, in any order, so
\texttt{mary smith} finds Smith, Mary. Punctuation does not matter:
\texttt{9045551234} and \texttt{904-555-1234} both work.

\subsection{The line under the boxes}

Something like: \textit{212 of 325 pilgrims, 3 of 25 Bambino, 0 on the
wait list, 7 drafts}.

Read that first number carefully. It counts \emph{walkers against the
cap}, not rows on the screen. It leaves out anyone who is not walking,
anyone you marked exempt, and anyone you added yourself, because none of
those count against the cap. So it will not match the number of lines
you can count in the table, and that is correct.

\subsection{The columns}

\begin{tabularx}{\linewidth}{@{}L{1.35in}Y@{}}
\toprule
\textbf{Party} & Everyone in the registration, with the email address
underneath. Press a name to open it. Press the email to write to them. \\
\textbf{State} & Where they live. Useful for the local versus
out-of-town question. \\
\textbf{Status} & Where they are up to. See the table below. \\
\textbf{Track} & \textit{full} is all three days. \textit{bambino} is
Sunday only. \\
\textbf{People} & How many are in the party, children included. \\
\textbf{Walkers} & How many of them actually walk. \\
\textbf{Created} & The day they signed up. \\
\textbf{Fees} & What the party owes in total. \\
\textbf{Paid} & How they paid, once they have. \\
\textbf{Waiver} & Whether the waiver is signed. \\
\textbf{Check-ins} & Ticks for the days of the walk they have arrived. \\
\textbf{Flags} & Small labels: exempt, assist, admin, non-walker, K/D
(Knight or Dame), vol (volunteer), refunded. \\
\bottomrule
\end{tabularx}

Press any column heading to sort by it. Press it again to sort the
other way.

\subsection{What the statuses mean}

This is the one piece of vocabulary worth learning. The words are the
program's, and they are short, so here is what each one actually means.

\closeup{img/cu-status.png}{The status column. On screen the
words are in capitals; that is only styling.}

\begin{tabularx}{\linewidth}{@{}L{1.35in}Y@{}}
\toprule
\textbf{draft} & They started the form and never finished. Nothing is
held for them. They are worth an email. \\[6pt]
\textbf{waitlist} & The track was full when they signed up. They have a
place in the queue and nothing else. \\[6pt]
\textbf{waiver} & \textbf{Needs to sign the waiver.} Their spots are
held. \\[6pt]
\textbf{payment} & \textbf{Needs to pay.} They have signed. Their spots
are held. \\[6pt]
\textbf{assistance} & They asked for help with the fee and are waiting
for your decision. \\[6pt]
\textbf{complete} & Signed and paid. Done. \\[6pt]
\textbf{cancelled} & Cancelled, by them or by you. Their spots went
back into the pool. You can put it back if it was a mistake. \\
\bottomrule
\end{tabularx}

\begin{note}{Good to know}
A spot is held for 48 hours from sign-up. If someone signs up and then
never pays, after two days their spot quietly goes back to the pool, so
an abandoned checkout does not keep a real pilgrim out. Their
registration stays, and they can still finish.
\end{note}

\subsection{Writing to several people at once}

The \btn{Copy 8 emails} button copies the email addresses of everyone
currently shown. Choose a segment first, say \btn{drafts}, press the
button, then paste into the To line of an email. That is how you chase
the people who never finished.

The number on the button can be smaller than the number of rows. A
sign-up that left only a telephone number has no address to copy, and
two people who typed the same address count once.

% ---------------------------------------------------------------- 4
\section{One registration}

Press any name on the roster. You get the whole registration: their
details on the left, and everything you can do about it on the right.

\shot{img/ov-registration.png}{The top of a registration. Their details
are a form you can change. Below it, and on a wide screen beside it,
come the cards shown on the next page.}

\subsection{Changing their details}

The left side is a form: contact details, the people in the party, what
each person is doing on each day, the parish, why they are walking, and
a box for your own notes that pilgrims never see.

Change what you need to, then press \btn{Save changes} at the bottom.

\begin{note}{Careful}
If you type something and then wander off to another screen without
saving, the program stops you and asks whether you really mean to leave.
Say no, then press \btn{Save changes}.
\end{note}

\subsection{The right-hand side}

\subsubsection{Status}

\closeup{img/reg-status.png}{The Status card.}

The status badge, the track, how they signed up, the fees, when they
signed up and when it last changed, and whether they are exempt from
the caps.

\subsubsection{Payment}

\closeup{img/reg-payment.png}{The Payment card, on a party that has
paid.}

What they paid, how, when, and any reference. If they still owe, you
get a \btn{Record this payment} form: choose \textit{check},
\textit{cash} or \textit{other}, put in the amount, a reference such as
a cheque number, and a note. That completes the registration.

If a payment exists, \btn{Mark this payment refunded} flags it. The
actual refund is done in Stripe; this only writes it on the record.

\subsubsection{Waiver}

\closeup{img/reg-waiver.png}{The Waiver card.}

Whether they signed, how, and when. \btn{They signed on paper} records
a paper waiver for somebody who could not sign online.

\subsubsection{Actions}

\closeup{img/reg-actions.png}{The Actions card. Which buttons appear
depends on where the registration has got to.}

Which buttons appear depends on the status.

\begin{tabularx}{\linewidth}{@{}L{1.85in}Y@{}}
\toprule
\btn{Offer them a spot} & On a wait-listed party. Moves them into the
queue to sign and pay. \\[6pt]
\btn{Approve assistance} & Completes the registration with no payment. \\[6pt]
\btn{Decline assistance} & Sends them to the payment step. \\[6pt]
\btn{Cancel this registration} & For the whole party. Their spots go
back. You are asked to confirm. \\[6pt]
\btn{Put it back} & On a cancelled one. Returns it to exactly the
status it had before. \\
\bottomrule
\end{tabularx}

Under the actions there is a box to pick an email template and press
\btn{Send this email again}.

\begin{note}{Good to know: the clock}
Times on these cards are shown in Greenwich time, which runs five hours
ahead of Florida in December. A sign-up at eight in the evening
therefore reads as the next day. Only the display is affected; the
program itself knows which day is which. This is on the list to put
right.
\end{note}

\subsubsection{Check-ins and History}

Check-ins show who arrived on which day, at what time, and who ticked
them off. \btn{Undo} takes one back.

\closeup{img/reg-history.png}{The History. Newest at the top.}

History is the whole story of the registration, newest first: every
change, who made it and when. Nothing is ever lost. This is what to
read when somebody asks what happened.

% ---------------------------------------------------------------- 5
\section{Adding a registration yourself}

For someone who telephoned, sent a cheque, or turned up on the day.

\shot{img/ov-add.png}{The Add screen. The same form a pilgrim fills in,
plus a card for what you took from them.}

\begin{enumerate}[leftmargin=*,itemsep=6pt]
\item Press \btn{Add a registration} on the roster.
\item Fill in the form the way the pilgrim would have.
\item On the right, under \btn{What you took from them}, tick
  \btn{Waiver signed on paper} if you have their signature on paper.
\item Now choose how they paid and put in the amount. \textbf{The
  payment boxes stay greyed out until the waiver tick is on}, because a
  pilgrim signs before they pay.
\item Press \btn{Add this registration}.
\end{enumerate}

\closeup{img/add-took.png}{What you took from them. The payment boxes
wake up once the waiver is ticked.}

The registration window and the caps do not apply here, because you
have looked and decided. These come out of the pool of late places the
organizers hold back.

% ---------------------------------------------------------------- 6
\section{The wait list}

When a track fills up, the public page stops taking sign-ups for it and
offers the wait list instead. A wait-listed party has a place in the
queue and nothing else: no waiver, no payment.

You decide who gets in. To offer somebody a spot:

\begin{enumerate}[leftmargin=*,itemsep=6pt]
\item Choose the \btn{wait list} segment on the roster.
\item Open the party you want.
\item Press \btn{Offer them a spot}.
\end{enumerate}

They move to the waiver step and their spots are held, whatever the
cap, because you have looked. Their number in the queue is shown beside
the status on the roster.

\begin{note}{Good to know}
Last year 50 to 60 people cancelled. Every cancellation frees spots,
and the wait list is where you fill them. It is worth checking the wait
list whenever a cancellation comes in.
\end{note}

% ---------------------------------------------------------------- 7
\section{The event weekend}

This is the part you will use on the three days themselves. There are
three ways a pilgrim can be checked in, and they all end up in the same
place.

\begin{enumerate}[leftmargin=*,itemsep=6pt]
\item \textbf{The pilgrim checks themselves in} from a link you email
  them that morning. This is the main way.
\item \textbf{A volunteer ticks them off} on a phone at the start line.
\item \textbf{You tick them off} from the Today screen.
\end{enumerate}

\subsection{The Today screen}

\shot{img/today.png}{The whole Today screen. Each part of it is shown
close up below.}

At the top: Friday, Saturday and Sunday. It opens on today if today is
one of the three, so most of the time you need not touch it.

\closeup{img/cu-count.png}{How the day is going, in one line.}

Under that: \textit{2 of 8 expected Friday checked in, 25\%} and
\textit{4 parties not here yet}. \textbf{Expected} means the people in
complete parties who are walking that day or coming to that day's
social. A Sunday-only family is not expected on Friday, and is not
counted.

\subsubsection{Emailing the day's check-in links}

\closeup{img/cu-send.png}{The send box. Here the button is greyed out
because today is not Friday.}

Each morning of the walk, press \btn{Email Friday's check-in links}
(the button names the day it will send). It goes to every complete
party who is expected that day, has not been checked in yet, and has
not already had that day's link. Press it again later in the morning to
catch the ones added since. Nobody gets two.

\begin{note}{Careful}
The button only works on the day itself. Looking at Saturday's tab on a
Friday, the button is greyed out and a second tick box appears: \btn{It
is not Saturday yet. Send anyway}. That is deliberate. It stops a
whole event's worth of emails going out a day early.
\end{note}

\subsubsection{The list}

\btn{Show} starts on \btn{someone missing}, which is what you want while
you are looking for people: a party stays on the list until the last of
them has arrived. The other choices are \btn{nobody here yet},
\btn{everyone here}, and \btn{everyone, any day}.

\closeup{img/cu-partyrow.png}{One party on the list, with a button for
the whole family and one for each person.}

Each row shows the party, whether they get a wristband, and buttons to
check them in. A circle means not arrived; a tick means arrived, with
who ticked them off and at what time. \textit{themselves} means the
pilgrim used their own link.

\btn{Check in everyone} does the whole party at once. The individual
names check in one person.

\subsection{What the pilgrim sees}

\phonepair{img/ov-pilgrim-checkin.png}{The pilgrim's own check-in link.
A party of three, so there is a box for each person.}{img/ov-checkin-app.png}{The
volunteers' app, on a phone at the start line.}

Somebody walking alone gets one button, \btn{I'm here}. A family gets a
box per person and one button, so a wife can check in before her
husband parks the car. Anyone who arrives later can use the same link
again.

Before the day, the link says when check-in opens. After the walk, it
says the Camino is over. A pilgrim who has not paid is told to see the
organizers at the start.

\begin{note}{Good to know}
A pilgrim can only tick people \emph{in}. They cannot undo anyone.
Taking a check-in back is yours and the volunteers' to do.
\end{note}

\subsection{The volunteers' phone app}

The volunteers use a different address, ending in \texttt{/checkin} instead of \texttt{/admin}.
They log in with the same password and put in their own first name, the
same way you do.

It is built for the beach. It works with no signal: the roster is kept
on the phone, and taps wait in a queue and go up when the phone finds a
signal again. The badge at the top says \textit{Synced} when everything
is up, \textit{3 to send} when taps are waiting, and \textit{Offline}
when the phone has no signal.

One tap on a name checks a person in and answers the wristband
question. \textbf{Wristband: give one} in green means they are paid and
signed. \textbf{No wristband:} in red gives the reason, for instance
\textit{unpaid} or \textit{waiver not signed}. The volunteer can still
check them in, and the reason stays on the record for you to sort out.

\begin{note}{Careful}
Have each volunteer open the app, log in, and put their name in
\emph{the day before}, at home, where there is a good signal. That way
the phone already holds the roster when they get to the beach.
\end{note}

\subsection{Counts}

The \btn{Counts} screen is where the actual head counts go: how many
were at Mass, on the walk, at the social, on the bus, each day. Type
the number, the time and a note, and press \btn{Save the counts}. These
sit beside the planned numbers in Reports, so next year you know how
many really came.

\btn{Time} and \btn{Note} are ordinary boxes: write whatever is useful,
such as \textit{just after eight}. A box left empty is not a zero, it
means nobody filled it in, and Reports treats it that way.

% ---------------------------------------------------------------- 8
\section{Reports and spreadsheets}

\shot{img/ov-reports.png}{The top of Reports. Everything here is worked
out from the roster, so it is always current.}

Reports has no buttons to press and changes nothing.

\closeup{img/reports-caps.png}{How full the event is, against every cap.}

It also shows:

\begin{itemize}[leftmargin=*,itemsep=5pt]
\item How full each track and each social is, against its cap.
\item How many are in each status, and on each track.
\item Payments, by method, with the amounts and any gifts. A refunded
  payment is still counted in this total, so for the exact picture give
  the treasurer the spreadsheet instead.
\item What is still owed, and how many have not signed.
\item Planned against actual, per day and per activity, with how many
  checked in.
\item How many Sunday walkers against the 350 the Shrine can feed.
\item Where people come from, which parishes, and sign-ups per day.
\end{itemize}

\closeup{img/reports-owing.png}{What is still outstanding.}

\subsection{The two spreadsheets}

At the top of Reports are two links. Press one and your computer
downloads a file you can open in Excel.

\begin{tabularx}{\linewidth}{@{}L{1.5in}Y@{}}
\toprule
\textbf{people.csv} & One line per person. Use this for name badges,
wristbands, bus lists and Mass counts. \\[6pt]
\textbf{regs.csv} & One line per registration, with the money and the
history. Use this for the treasurer. \\
\bottomrule
\end{tabularx}

% ---------------------------------------------------------------- 9
\section{Settings}

\shot{img/ov-settings.png}{The top of Settings. Change what you need,
then press \textnormal{\textbf{Save the settings}} at the very bottom.}

This is where the event itself is described. Most of it is set once,
before registration opens, and then left alone.

\closeup{img/settings-event.png}{The event itself: its name, the clock
setting, the three dates and the public address.}

\closeup{img/settings-fees.png}{The fees. In dollars, not cents.}

\closeup{img/settings-window.png}{The window: when sign-ups open, when
they close, and the last moment a pilgrim may change their own.}

\begin{tabularx}{\linewidth}{@{}L{1.85in}Y@{}}
\toprule
\textbf{Name} & What the event is called. \\[5pt]
\textbf{Hours from UTC} & Leave it at $-5$. It tells the program what
day it is in Florida, so check-in opens at the right hour. \\[5pt]
\textbf{Days, one per line} & The three dates of the walk, earliest
first, written year first: \texttt{2026-12-04}. \\[5pt]
\textbf{Public URL} & The address pilgrims use. \\[5pt]
\textbf{Fees} & In dollars. Per person, children included. \\[5pt]
\textbf{Caps} & How many may sign up for each track and each social,
and how many late places you hold back. The Shrine number is reported
only, never enforced. \\[5pt]
\textbf{Window} & \textbf{Opens} is when the public page starts taking
sign-ups. \textbf{Closes} is when it stops. \textbf{Changes close} is
the last moment a pilgrim may edit their own registration. These three
show the time on \emph{your own computer}. If you are setting them up
from another time zone, allow for it. \\[5pt]
\textbf{Hold hours} & How long a spot is held for somebody who has not
finished paying. 48 is the normal answer. \\[5pt]
\textbf{Organizations} & The parishes and groups the sign-up form
suggests. One per line. \\[5pt]
\textbf{Providers} & Leave alone. See the box below. \\
\bottomrule
\end{tabularx}

\begin{note}[warnbg]{Careful: rehearsal mode}
While \textbf{Mode} says \textit{stub}, the program is in rehearsal: no
email is really sent, and payments complete themselves without any
money changing hands. You will see the word \textit{stub} in the Paid
column and a line saying \textit{Rehearsal: no email leaves the ship}.
That is correct for testing. Switching it to \textit{live} is not
something to do on your own; ask whoever set the program up.
\end{note}

% ---------------------------------------------------------------- 10
\section{The emails}

\shot{img/ov-emails.png}{The Emails screen. One card per email the
program sends.}

Each email has a subject and a body you can edit. Press \btn{Save this
email} on that card when you are done.

The curly-bracket words, such as \verb|{{first}}| and \verb|{{link}}|,
are filled in by the program when it sends: the pilgrim's first name,
their private link, the total, and so on. The card lists which ones
that email uses.

\begin{note}{Careful}
Do not delete a curly-bracket word. If you do, the program refuses the
save and says \textit{Put these back before saving}, and shows you
which. Put it back and save again.

The refusal covers the whole card. If you improved the subject and then
lost a placeholder in the body, both go back to how they were, so check
the subject again before you save a second time.
\end{note}

You may move a curly-bracket word, put words around it, or use it
twice. You simply may not lose one.

Every other piece of writing that a pilgrim reads is edited on the
public page itself, not here. See the next section.

% ---------------------------------------------------------------- 11
\section{Changing the words pilgrims read}

Everything on the sign-up page, every heading, every label, every bit of
help text, can be rewritten by you, in place, on the page itself.

\begin{enumerate}[leftmargin=*,itemsep=6pt]
\item Open the public page, \texttt{forgave-infect.nisfeb.com/apps/register/}, while you are logged
  in.
\item Press \btn{Edit text} at the top.
\item Click any piece of writing and type over it.
\item Click away. A tick appears when it is saved.
\end{enumerate}

A row of buttons appears, \btn{Show all steps}, with one for each page a
pilgrim might see: the landing page, the form, the waiver step, the
payment step, the wait list, and so on. Press one to see that page,
filled with a made-up example, so you can edit its words without having
to sign up yourself.

\begin{note}{Careful}
While \btn{Edit text} is on, nothing else on the page works: buttons do
nothing and boxes cannot be typed in. That is on purpose, so you cannot
accidentally sign yourself up while editing. Press \btn{Editing text}
again to turn it off.
\end{note}

% ---------------------------------------------------------------- 12
\section{Backup}

\shot{img/ov-backup.png}{Backup. The downloads at the top, the restore
below.}

\subsection{Taking a backup}

Press \btn{JSON bundle} or \btn{jam} and save the file somewhere safe:
your own computer, and ideally a second place. Do this weekly once
registration opens, and every evening during the walk.

The \textbf{jam} is the emergency copy and the exact one. The
\textbf{JSON bundle} holds the same information in a form a person can
read. Both leave out passwords and keys.

\subsection{Putting a backup back}

\begin{enumerate}[leftmargin=*,itemsep=6pt]
\item Press \btn{Choose file} and pick the backup.
\item Press \btn{Read the file}. The program tells you what is in it:
  how many registrations, the dates they cover, and how many of your
  current ones it would overwrite. It changes nothing yet.
\item If it looks right, press \btn{Restore} and confirm.
\end{enumerate}

\begin{note}[warnbg]{Careful}
The tick box \btn{Delete every registration the file does not name} is
for a true restore onto an empty system. Leave it \emph{off} unless you
are certain. A restore also stamps every registration as changed now,
which restarts the 48-hour hold on anyone who has not finished paying.
\end{note}

Do the three steps in one sitting. If you wander off to another screen,
or somebody signs up while you are deciding, the program asks you to
press \btn{Read the file} again before it will restore. That is on
purpose: it will only restore onto the situation it showed you.

% ---------------------------------------------------------------- 13
\section{If something looks wrong}

\begin{tabularx}{\linewidth}{@{}L{1.95in}Y@{}}
\toprule
\textbf{It asks me to log in again} & Normal after about a week, or on a
new computer. Type the password again. \\[7pt]
\textbf{The page says the ship would not take that change} & Something
was refused. Look at the registration again: it will show what the
program actually has. Then try once more. \\[7pt]
\textbf{A name I know is registered does not come up in Search} & Check
the Segment box. If it is on \btn{registrations}, a cancelled or
abandoned one will not show. Choose \btn{everything}. \\[7pt]
\textbf{The page is slow} & Give it a few seconds. Everything you press
shows immediately and settles afterwards; you do not need to press
twice. \\[7pt]
\textbf{A pilgrim says their link does not work} & Open their
registration and check the status. If they are cancelled, the link says
so. Send them the right email again from the box under Actions. \\[7pt]
\textbf{I cancelled the wrong party} & Open it and press \btn{Put it
back}. It returns to exactly where it was. \\[7pt]
\textbf{Two of us edited the same registration} & The last save wins.
Read the History at the bottom: it shows both changes, with names and
times. \\
\bottomrule
\end{tabularx}

\begin{note}{Good to know}
Nothing you do here is truly lost. Every registration keeps its whole
history, and cancellations can be put back. If you are unsure, look at
the History before you worry.
\end{note}

% ---------------------------------------------------------------- 14
\section{Words this program uses}

\begin{tabularx}{\linewidth}{@{}L{1.5in}Y@{}}
\toprule
\textbf{Backoffice} & Your side of the program, at the /admin address. \\
\textbf{Bambino} & The Sunday-only track, the last 2.5 miles. \\
\textbf{Cap} & The most people allowed on a track or at a social. \\
\textbf{Complete} & Signed and paid. Nothing outstanding. \\
\textbf{Draft} & A sign-up somebody started and never finished. \\
\textbf{Exempt} & Marked as not counting against the caps. Used for
local Order of Malta members and volunteers. \\
\textbf{Full track} & All three days of the walk. \\
\textbf{Hold} & The 48 hours a spot is kept for somebody who has signed
up but not finished paying. \\
\textbf{Party} & One registration, which may be one person or a whole
family. \\
\textbf{Rehearsal (stub)} & Practice mode: no real emails, no real
money. \\
\textbf{The ship} & The computer this program runs on. It shows up in a
few messages. \\
\textbf{Track} & Which walk they are on: full or Bambino. \\
\textbf{Waiver} & The liability form every pilgrim signs. \\
\textbf{Wristband} & What a checked-in pilgrim is given. Green means
give one. \\
\bottomrule
\end{tabularx}

% ---------------------------------------------------------------- 15
\newpage
\section*{The one-page reminder}
\addcontentsline{toc}{section}{The one-page reminder}

\textbf{Your address:} \texttt{forgave-infect.nisfeb.com/apps/register/admin}

\textbf{Every day of registration}

\begin{enumerate}[leftmargin=*,itemsep=4pt]
\item Open the roster. Read the line under the boxes for the totals.
\item Segment \btn{pending}: who still owes a waiver or a payment.
\item Segment \btn{financial assistance}: decisions waiting for you.
\item Segment \btn{drafts}: press \btn{Copy emails} and chase them.
\item Segment \btn{wait list} when somebody cancels.
\end{enumerate}

\textbf{Every week}

\begin{enumerate}[leftmargin=*,itemsep=4pt]
\item \btn{Backup}, press \btn{jam}, save the file somewhere safe.
\item \btn{Reports} for the totals and the money.
\end{enumerate}

\textbf{Each morning of the walk}

\begin{enumerate}[leftmargin=*,itemsep=4pt]
\item \btn{Today}. Check it is on the right day.
\item Press the \btn{Email \ldots\ check-in links} button. It names the
  day, so on Friday it reads \btn{Email Friday's check-in links}.
\item Press it again mid-morning for anyone added since.
\item Watch the percentage. Use \btn{Show: someone missing} to find who
  is not here.
\item Check people in by hand when their phone will not cooperate.
\end{enumerate}

\textbf{Each evening of the walk}

\begin{enumerate}[leftmargin=*,itemsep=4pt]
\item \btn{Counts}: put in the day's real head counts.
\item \btn{Backup}: take a copy.
\end{enumerate}

\vfill
{\small This manual describes Register version 13.\par}

\end{document}
